% !TEX root = ../main.tex
\section{Background and Related Work}
\label{sec:drl_robotics}

This chapter develops the theory the thesis relies on and situates the work
within the surrounding literature. It begins with the formal definitions of
reinforcement learning, first for fully observable problems and then for the
partially observable case that motivates the recurrent policy used here. It then
describes the two algorithms used in the experiments, Proximal Policy
Optimization and Soft Actor-Critic, before describing the two techniques the
thesis compares: potential-based reward shaping, and structured training through
curriculum learning and asymmetric self-play. The chapter ends with the
mechanics and geometry of planar pushing, which explain why the task is
multi-objective and why its two objectives cannot be optimised in
isolation.

\subsection{Reinforcement Learning}

Reinforcement Learning (RL) is the branch of machine learning concerned with how
an agent should act in an environment so as to maximise a cumulative scalar
reward. Unlike supervised learning, the agent is never told the correct action;
it receives only an evaluative signal and must discover, through trial and error,
which behaviours lead to high long-term return. The agent observes the state of
the environment, selects an action, and receives a reward together with a new
state; repeating this loop produces a trajectory whose total reward the agent
learns to improve. This interaction is formalised through the Markov Decision
Process.

\subsubsection{Markov Decision Processes}

A Markov Decision Process (MDP) is the standard model of a fully observable
sequential decision problem. Following \cite{Ng1999PolicyInvariance,
Sutton1998}, it is a tuple $\mathcal{M} = (\mathcal{S}, \mathcal{A}, P, R,
\gamma)$, where:
\begin{itemize}
  \item $\mathcal{S}$ is the \textbf{state space}, the set of all configurations
        the environment can occupy. At time step $t$ the agent perceives a state
        $s_t \in \mathcal{S}$.
  \item $\mathcal{A}$ is the \textbf{action space}, the set of actions available
        to the agent; in robotics this may be a set of joint commands or, as in
        this thesis, a set of parameterised motion primitives.
  \item $P(s_{t+1} \mid s_t, a_t)$ is the \textbf{transition function}, giving the
        probability of reaching $s_{t+1}$ after taking action $a_t$ in state
        $s_t$. It encodes the environment dynamics.
  \item $R(s_t, a_t, s_{t+1})$ is the \textbf{reward function}, the scalar
        feedback for a transition.
  \item $\gamma \in [0, 1]$ is the \textbf{discount factor}, which sets the
        relative weight of immediate against future reward.
\end{itemize}

The defining property is the \emph{Markov} assumption: the transition and reward
depend only on the current state and action, not on the history that preceded
them. Formally, $P(s_{t+1} \mid s_t, a_t) = P(s_{t+1} \mid s_0, a_0, \dots, s_t,
a_t)$. The state is therefore a sufficient statistic of the past, which is what
makes value functions and the Bellman recursion well defined.

The agent acts according to a \textbf{policy} $\pi(a \mid s)$, a mapping from
states to a distribution over actions. Its objective is to maximise the
\textbf{expected return}, the expected discounted sum of future rewards:
\begin{equation}
    J(\pi) = \mathbb{E}_{\pi}\!\left[ \sum_{t=0}^{\infty} \gamma^t R(s_t, a_t, s_{t+1}) \right].
    \label{eq:return}
\end{equation}
The discount factor $\gamma$ has a dual role. Mathematically it keeps the
infinite sum bounded; behaviourally it interpolates between a myopic agent
($\gamma \to 0$), which optimises only the immediate reward, and a far-sighted
agent ($\gamma \to 1$), which weighs distant rewards almost as heavily as
immediate ones. In episodic tasks such as the pushing problem studied here, the
horizon is finite and the undiscounted case $\gamma = 1$ is admissible provided
every trajectory eventually terminates, a subtlety that becomes important for
reward shaping in Section~\ref{sec:pbrs}.

\paragraph{Value functions.}
To reason about long-term consequences, RL introduces two value functions. The
\textbf{state-value function} $V^{\pi}(s)$ is the expected return obtained by
starting in state $s$ and following $\pi$ thereafter,
\begin{equation}
    V^{\pi}(s) = \mathbb{E}_{\pi}\!\left[ \sum_{t=0}^{\infty} \gamma^t R(s_t, a_t, s_{t+1}) \;\middle|\; s_0 = s \right],
\end{equation}
while the \textbf{action-value function} $Q^{\pi}(s, a)$ conditions additionally
on taking action $a$ first,
\begin{equation}
    Q^{\pi}(s, a) = \mathbb{E}_{\pi}\!\left[ \sum_{t=0}^{\infty} \gamma^t R(s_t, a_t, s_{t+1}) \;\middle|\; s_0 = s,\, a_0 = a \right].
\end{equation}
The two are linked by $V^{\pi}(s) = \mathbb{E}_{a \sim \pi(\cdot \mid s)}[Q^{\pi}(s,
a)]$, and their difference defines the \textbf{advantage} $A^{\pi}(s, a) =
Q^{\pi}(s, a) - V^{\pi}(s)$, which measures how much better action $a$ is than
the policy's average behaviour in state $s$. The advantage is central to the
policy-gradient methods used in this thesis.

\paragraph{Bellman equations.}
Value functions satisfy a recursive consistency condition. For the action-value
function,
\begin{equation}
    Q^{\pi}(s, a) = \mathbb{E}_{s' \sim P(\cdot \mid s, a)}\!\left[ R(s, a, s') + \gamma\, V^{\pi}(s') \right],
    \label{eq:bellman}
\end{equation}
which decomposes the value of a state-action pair into the immediate reward and
the discounted value of the successor state. The optimal value functions
$V^{*}(s) = \sup_{\pi} V^{\pi}(s)$ and $Q^{*}(s, a) = \sup_{\pi} Q^{\pi}(s, a)$
induce the optimal policy $\pi^{*}(s) = \arg\max_{a} Q^{*}(s, a)$. Whether solved
by dynamic programming, value iteration, or sampled approximation, essentially
all RL algorithms can be seen as ways of estimating or improving the quantities
in Equation~\ref{eq:bellman}.

\paragraph{Exploration versus exploitation.}
Because the agent learns from its own experience, it faces a fundamental
trade-off. \emph{Exploitation} chooses the action currently believed to be best,
which secures known reward; \emph{exploration} tries alternatives to improve
those beliefs, at the risk of short-term loss for long-term gain. An agent that
only exploits may converge to a suboptimal behaviour, while one that only
explores never applies what it has learned. Mechanisms such as $\epsilon$-greedy
action selection and entropy regularisation, the latter central to Soft
Actor-Critic (Section~\ref{sec:sac}), exist to manage this trade-off. In
sparse-reward manipulation the trade-off is difficult, because informative reward
is encountered so rarely that undirected exploration is unlikely to reach it, a
difficulty that both reward shaping and automatic curricula attempt to address.

\subsubsection{Partial Observability and the POMDP}
\label{sec:pomdp}

The MDP assumes the agent perceives the full state. Real and simulated robots
rarely enjoy this: sensors are noisy, occlusions hide part of the scene, and
quantities such as contact forces or friction coefficients are simply not
observable. Such problems are modelled as a Partially Observable Markov Decision
Process (POMDP), a tuple $(\mathcal{S}, \mathcal{A}, P, R, \Omega, O, \gamma)$
that augments the MDP with an \textbf{observation space} $\Omega$ and an
\textbf{observation function} $O(o \mid s', a)$ giving the probability of
observation $o$ after a transition to $s'$. The agent no longer sees $s_t$; it
sees only $o_t$, and a single observation is generally not a sufficient statistic
of the state.

Because the observation breaks the Markov property, an optimal policy cannot in
general be a function of the current observation alone. The classical solution is
to act on a \textbf{belief state}, the posterior $b_t(s) = P(s_t = s \mid o_0,
a_0, \dots, o_t)$ over states given the entire interaction history, which is
itself Markovian; maintaining the exact belief, however, is intractable for
continuous state spaces. A practical alternative, and the one adopted in this
thesis, is to have the policy summarise the history with a recurrent network. A
Long Short-Term Memory (LSTM) network maintains a hidden state that is updated at
each step, which provides the policy with a learned, compressed approximation of
the belief. In the pushing task the current observation reports object and goal
poses, but not the underlying contact state or the outcome of prior pushes. The
problem is therefore formally a POMDP. Conditioning the policy on the sequence of
pushes through an LSTM lets it combine information across the episode, rather than
respond to a single frame. For this reason every policy compared in this thesis
is recurrent.

\subsubsection{Model-Free and Model-Based Reinforcement Learning}

RL algorithms divide into two families according to whether they build an
explicit model of the environment dynamics $P(s_{t+1} \mid s_t, a_t)$
\cite{Yang2023SimtoRealMA}. \textbf{Model-free} methods learn a policy or value
function directly from sampled interaction, without ever estimating the
transition function. They subdivide further:
\begin{itemize}
  \item \textbf{Value-based} methods, such as Deep Q-Networks \cite{DQN}, learn
        an action-value function and act greedily with respect to it, obeying the
        recursion $Q(s, a) = \mathbb{E}[R(s, a, s') + \gamma \max_{a'} Q(s',
        a')]$.
  \item \textbf{Policy-based} methods, such as PPO
        \cite{Schulman2017ProximalPO}, parameterise the policy directly and
        ascend the gradient of the expected return in Equation~\ref{eq:return}.
  \item \textbf{Actor-critic} methods, such as Soft Actor-Critic
        \cite{Haarnoja2018SoftAO}, combine the two, maintaining an actor that
        represents the policy and a critic that estimates value to reduce the
        variance of the policy gradient.
\end{itemize}
Model-free methods make no assumptions about the dynamics and are robust to
complex, contact-rich physics, which makes them the default choice in
manipulation; their weakness is high sample complexity, since every unit of
information must be extracted from interaction.

\textbf{Model-based} methods instead learn an approximate dynamics model
$\hat{P}(s_{t+1} \mid s_t, a_t)$ and use it to plan or to generate synthetic
experience, typically by minimising a prediction error such as
\begin{equation}
    L_{\text{model}} = \mathbb{E}_{\mathcal{D}}\!\left[ \left\lVert P(s_{t+1} \mid s_t, a_t) - \hat{P}(s_{t+1} \mid s_t, a_t) \right\rVert^2 \right],
\end{equation}
over a dataset $\mathcal{D}$ of observed transitions. They can be far more
sample-efficient, but their performance is bounded by the fidelity of the learned
model, and small errors compound over long rollouts. The choice between the two
families is therefore a choice between robustness and sample efficiency. This
thesis works entirely in the model-free, on-policy regime and asks a
complementary question: rather than learning a dynamics model to save samples,
can a principled reward do so instead, by making each sampled transition more
informative? Section~\ref{sec:pbrs} takes up this question.

\subsection{Proximal Policy Optimization}
\label{sec:ppo}

Proximal Policy Optimization (PPO) \cite{Schulman2017ProximalPO} is the on-policy,
model-free algorithm used to train every policy in this thesis. It belongs to the
policy-gradient family, which optimises a parameterised policy $\pi_{\theta}(a
\mid s)$ by ascending an estimate of the gradient of the expected return. The
basic policy-gradient estimator,
\begin{equation}
    \nabla_{\theta} J(\pi_{\theta}) = \mathbb{E}_{\tau \sim \pi_{\theta}}\!\left[ \nabla_{\theta} \log \pi_{\theta}(a_t \mid s_t)\, \hat{A}_t \right],
\end{equation}
scales the score function of each action by its estimated advantage $\hat{A}_t$,
so that actions better than average are made more likely. In practice the
advantage is computed with Generalized Advantage Estimation (GAE), which trades
bias against variance by exponentially weighting multi-step temporal-difference
errors; GAE reappears in Chapter~\ref{section:methods}, where its decomposition
explains how reward shaping interacts with a stale value estimate.

PPO improves on the earlier Trust Region Policy Optimization
\cite{Schulman2015TrustRP}, which constrained each update to a trust region
through a costly second-order optimisation. PPO achieves a similar effect with a
first-order \textbf{clipped surrogate objective}. Writing the probability ratio
between the new and old policies as $r_t(\theta) = \pi_{\theta}(a_t \mid s_t) /
\pi_{\theta_{\text{old}}}(a_t \mid s_t)$, the objective is
\begin{equation}
    L^{\text{CLIP}}(\theta) = \mathbb{E}_t\!\left[ \min\!\left( r_t(\theta)\, \hat{A}_t,\; \operatorname{clip}\!\big(r_t(\theta), 1 - \epsilon, 1 + \epsilon\big)\, \hat{A}_t \right) \right].
    \label{eq:ppo_clip}
\end{equation}
The clipping removes the incentive to move the policy ratio beyond $[1 - \epsilon,
1 + \epsilon]$, which keeps each update inside an implicit trust region and
prevents the large, destabilising steps that occur in unconstrained policy
gradients. This same clipping mechanism is later reused in the behavioural-cloning
loss of asymmetric self-play (Section~\ref{sec:asp}), where it has an unintended
consequence that this thesis diagnoses. The full PPO objective adds a value-function
loss $L^{\text{VF}}(\phi) = \mathbb{E}_t[(V_{\phi}(s_t) - \hat{R}_t)^2]$ and an
entropy bonus $L^{\text{ENT}} = \mathbb{E}_t[\mathcal{H}(\pi_{\theta}(\cdot \mid
s_t))]$ that encourages exploration:
\begin{equation}
    L(\theta) = L^{\text{CLIP}}(\theta) - c_1 L^{\text{VF}}(\phi) + c_2 L^{\text{ENT}},
\end{equation}
with coefficients $c_1$ and $c_2$ balancing the three terms.

\subsection{Soft Actor-Critic}
\label{sec:sac}

Soft Actor-Critic (SAC) \cite{Haarnoja2018SoftAO} is an off-policy actor-critic
algorithm built on the maximum-entropy framework. Where standard RL maximises
return alone, SAC augments the objective with the entropy of the policy,
\begin{equation}
    J(\pi) = \sum_{t=0}^{T} \mathbb{E}_{(s_t, a_t) \sim \rho_{\pi}}\!\left[ R(s_t, a_t) + \alpha\, \mathcal{H}\big(\pi(\cdot \mid s_t)\big) \right],
    \label{eq:sac}
\end{equation}
where $\mathcal{H}(\pi(\cdot \mid s_t)) = \mathbb{E}[-\log \pi(\cdot \mid s_t)]$
is the policy entropy and the temperature $\alpha$ sets the weight of exploration
against reward. Maximising entropy encourages the policy to remain as stochastic
as the reward permits. This improves exploration and robustness, and discourages
premature convergence to a single mode, an advantage in sparse-reward or
multimodal problems. Being off-policy, SAC reuses past experience from a replay
buffer, and is typically far more sample-efficient than on-policy methods. SAC is
relevant here as the basis of an off-policy comparison baseline. Combined with
Hindsight Experience Replay \cite{Andrychowicz2017}, which relabels failed
trajectories with the goals they happened to achieve, it provides a natural
goal-conditioned learner against which the on-policy results can be calibrated, as
described in Chapter~\ref{sec:experiments}.

\subsection{Potential-Based Reward Shaping}
\label{sec:pbrs}

Reward shaping supplies an agent with additional reward beyond the sparse task
signal, to guide learning through the temporal credit-assignment problem
\cite{Grzes2017RewardShaping}. Its danger is equally well known: an ill-chosen
shaping reward can change the optimal policy, so that the agent learns to exploit
the shaping rather than to solve the task. \cite{Ng1999PolicyInvariance} give the
canonical illustration, a bicycle agent rewarded for progress toward a goal that
learns to ride in circles near the start, collecting the progress reward
repeatedly without ever advancing.

Potential-Based Reward Shaping (PBRS) is the principled resolution of this
danger. \cite{Ng1999PolicyInvariance} prove that a shaping reward of the form
\begin{equation}
    F(s, a, s') = \gamma\, \Phi(s') - \Phi(s),
    \label{eq:pbrs}
\end{equation}
defined as the discounted difference of a \textbf{potential function}
$\Phi: \mathcal{S} \to \mathbb{R}$ over states, leaves the optimal policy of the
MDP unchanged for \emph{any} choice of $\Phi$. The intuition is that the shaping
term telescopes along a trajectory: summed over an episode it collapses to a
boundary term that depends only on the first and last states, so it cannot create
the reward-collecting cycles that non-potential shaping permits. Formally, adding
Equation~\ref{eq:pbrs} to the task reward shifts every action-value by the same
state-dependent constant, $Q'(s, a) = Q(s, a) - \Phi(s)$, and since the shift is
independent of the action it does not alter the arg-max that defines the optimal
policy. This decoupling is the property the thesis exploits: it permits a dense,
informative reward to be designed purely to accelerate learning, with a guarantee
that the task being solved is unchanged.

Subsequent work extended the framework in directions this thesis uses.
\cite{Devlin2012DynamicPBRS} generalised the potential to depend on time as well
as state, $F(s, t, s', t') = \gamma \Phi(s', t') - \Phi(s, t)$. This preserves
policy invariance while letting the shaping adapt during learning, and so licences
shaping that changes across a curriculum. \cite{Harutyunyan2015} showed that an
arbitrary reward function can be expressed as potential-based advice, widening the
class of admissible potentials. Crucial for the present setting is the episodic
case. \cite{GrzesKudenko2009} and \cite{Grzes2017RewardShaping} analysed PBRS in
the finite-horizon setting and clarified when the undiscounted case $\gamma = 1$ is
safe. Provided every trajectory terminates and terminal states are handled
consistently, the shaping sum telescopes to $\Phi(s_T) - \Phi(s_0)$ and remains
bounded. This result justifies the use of $\gamma_{\text{shaping}} = 1$ in the
potentials constructed in Chapter~\ref{section:methods}. The theorem and its
episodic refinement are the foundation of the reward design in this thesis. The
specific potentials, which combine bounded exponential functions of position and
orientation error, are introduced there rather than here, since they are a
contribution of the work rather than prior art.

\subsection{Curriculum Learning and Automatic Curricula}
\label{sec:curriculum}

An orthogonal approach to sample inefficiency is to change not the reward but
the order in which the agent encounters tasks. Curriculum learning organises
training from easy to hard, so that competence acquired on simple problems
supports progress on difficult ones \cite{Narvekar2020}. The survey of
\cite{Narvekar2020} provides a taxonomy that distinguishes, among other axes,
\emph{forced} curricula, in which a designer fixes the sequence of tasks or
objectives in advance, from \emph{automatic} curricula, in which the sequence is
derived from the learner's own progress. Both appear in this thesis: a forced
position-then-rotation curriculum, and an automatic curriculum generated by
self-play.

Several forced and semi-automatic schemes inform the design. Reverse curriculum
generation \cite{Florensa2017ReverseCurriculum} starts the agent near the goal
and expands the start distribution outward as competence increases, so that the
agent always trains at the boundary of what it can already achieve. Precision-based
curricula such as Continuous Curriculum Learning \cite{Luo2020AcceleratingRL}
progressively tighten the success tolerance, which is a close analogue of staging
a pushing task from coarse positioning to fine orientation control. The survey of
\cite{Portelas2020ACLSurvey} organises automatic curriculum methods into families
that include self-paced learning, teacher-student setups, and goal-generation, and
notes the recurring difficulty that an automatic curriculum can propose goals
faster than the learner masters them. That failure mode is the one the self-play
experiments in this thesis encounter.

\subsection{Asymmetric Self-Play}
\label{sec:asp}

Asymmetric Self-Play (ASP) is the automatic curriculum at the centre of the
thesis. Introduced by \cite{Sukhbaatar2018IntrinsicMotivation}, it pits two
policies embodied in the same agent against each other: Alice proposes a task by
acting in the environment, and Bob must then solve it. Because Alice
demonstrates each task by performing it, every proposed goal is achievable by
construction, and the competition between the two produces a curriculum of
increasing difficulty with no external reward or human specification. Sukhbaatar
et al. define a self-regulating reward structure,
\begin{align}
    R_B &= -\gamma\, t_B, \label{eq:sukh_bob}\\
    R_A &= \gamma \max(0,\, t_B - t_A), \label{eq:sukh_alice}
\end{align}
where $t_A$ is the time Alice takes to set the task, $t_B$ the time Bob takes to
solve it, and $\gamma$ a scaling coefficient. Alice is rewarded when Bob is slow
but penalised for her own effort, so her optimal behaviour is to find the
\emph{simplest} tasks that Bob cannot yet solve, placing each new task just beyond
his current ability. This time-based reward is one of the two Alice reward
formulations the thesis evaluates.

\cite{OpenAI2021AsymmetricSF} scaled ASP to robotic manipulation and introduced
the second formulation. They model the problem as a goal-augmented MDP
$\mathcal{M} = \langle \mathcal{S}, \mathcal{A}, \mathcal{P}, \mathcal{R},
\mathcal{G} \rangle$ with an explicit goal space $\mathcal{G}$, run Alice to
produce a trajectory whose final state $s_T$ becomes the goal for Bob, and train
Bob from the same initial state on the resulting goal distribution $\mathcal{G}(g
\mid s_0)$. Their Alice reward is \emph{outcome-based} rather than time-based:
Alice receives a reward when Bob fails to reach the goal, a shaping bonus for
proposing a valid goal, and a penalty for placing objects out of bounds, while
Bob receives a sparse reward for reaching the goal. Both the outcome-based and
the time-based Alice rewards are compared in this thesis.

A distinctive feature of the manipulation formulation is that Alice's own
trajectory is itself a demonstration of the goal she sets, which
\cite{OpenAI2021AsymmetricSF} use through Alice Behavioral Cloning (ABC). Bob
is trained on a combined objective $\mathcal{L} = \mathcal{L}_{\text{RL}} + \beta
\mathcal{L}_{\text{abc}}$, where the behavioural-cloning term imitates Alice's
relabelled trajectory on goals Bob failed to solve. To prevent the imitation term
from producing destabilising policy updates, they apply the PPO clipping mechanism
of Equation~\ref{eq:ppo_clip}, with the advantage fixed to one:
\begin{equation}
    \mathcal{L}_{\text{abc}} = -\,\mathbb{E}_{(s_t, g_t, a_t) \in \mathcal{D}_{\text{BC}}}\!\left[ \operatorname{clip}\!\left( \frac{\pi_B(a_t \mid s_t, g_t; \theta)}{\pi_B(a_t \mid s_t, g_t; \theta_{\text{old}})},\, 1 - \epsilon,\, 1 + \epsilon \right) \right].
    \label{eq:abc}
\end{equation}
This clipped behavioural-cloning loss behaves as intended only when Alice's
demonstrated actions lie close to Bob's current policy. When the two diverge
strongly, as they do under the discrete, contact-rich action space of this thesis,
the clip saturates and the imitation gradient vanishes; this interaction is
analysed in the results. Behavioural cloning more broadly provides the wider
context for ABC. This includes cloning from observation \cite{Torabi2018BCO},
implicit energy-based cloning \cite{Florence2022ImplicitBC}, and
demonstration-boosted RL \cite{Hester2017DeepQF, Silver2018ResidualPolicy}, all
methods for adding a demonstration signal to an RL objective.

The broader set of automatic curricula based on adversarial goal generation
places ASP among several related methods. PAIRED \cite{Dennis2020PAIRED} generates
environments that maximise the regret between a protagonist and an antagonist.
AMIGo \cite{Campero2021AMIGo} trains a teacher to propose goals a student finds
challenging but achievable. Adversarial Intrinsic Motivation
\cite{Durugkar2021AIM} directs exploration with a Wasserstein distance to the goal
distribution. All share ASP's central premise, that a well-tuned adversary
produces a useful curriculum, and all share its central risk, that the adversary
may generate tasks the learner cannot use. The stability of such coupled
optimisation is itself a subject of study. \cite{Letcher2019GameMechanics}
decompose the dynamics of differentiable games into potential and Hamiltonian
components, and identify when adversarial training converges or cycles.
\cite{Berner2019Dota2} demonstrate that competitive self-play can reach very high
performance, but only at very large scale. These two observations define the
scope of the present study: the self-play results here are obtained under a
single-Graphical-Processing-Unit budget, and are interpreted accordingly.

The goal criterion also separates this thesis from the settings in which ASP
succeeded, and the difference determines what the method must achieve. In the
original formulation of \cite{Sukhbaatar2018IntrinsicMotivation}, a goal is a
single state Bob must reach, often in a domain where position dominates and
orientation is absent or loosely constrained. In the manipulation setting of
\cite{OpenAI2021AsymmetricSF}, a goal is the final pose of Alice's trajectory,
scored within a tolerance on each component. The task studied here instead applies
a strict combined gate. A goal counts as reached only when position and orientation
both lie within tolerance at the same time. The contact mechanics of
Section~\ref{sec:pushing} couple the two, so neither can be optimised alone. The
goal is therefore genuinely multi-objective in a sense the prior settings did not
require. This is a difference in regime, not a defect in the earlier work, and it
defines one of the axes along which this thesis explains why a proven method
underperforms here.

\subsection{Goal-Conditioned and Hierarchical Reinforcement Learning}
\label{sec:gcrl}

ASP produces a goal-conditioned policy, and the way a goal is represented
influences what the policy can learn. Goal-conditioned RL augments the state with
a goal and learns a single policy $\pi(a \mid s, g)$ that generalises across
goals; Hindsight Experience Replay \cite{Andrychowicz2017} makes this practical
under sparse reward by relabelling achieved states as goals. Sukhbaatar et al. in
their goal-embedding work \cite{Sukhbaatar2018GoalEmbeddings} showed that
compressing goals into a learned latent through a goal encoder can improve
hierarchical learning, an idea the thesis adopts in the form of the GoalEncoder
used by Bob. The effect of such a compression can be positive or negative, and the
information-bottleneck literature explains why. \cite{Tishby2015InfoBottleneck}
formalise the trade-off between compression and predictive sufficiency. InfoBot
\cite{Goyal2019InfoBot} applies it in RL, and shows that a bottleneck can either
improve generalisation by discarding distractors or reduce precision by discarding
task-relevant detail. Whether the goal encoder improves or reduces performance in
the pushing task is one of the questions the ablations address.

Hierarchical RL provides further context for the two-level Alice-Bob structure.
Methods such as HIRO \cite{Nachum2018DataEfficientHR} and FeUdal Networks
\cite{Vezhnevets2017FeUdal} decompose control into a high-level policy that sets
sub-goals and a low-level policy that achieves them. The Option-Critic architecture
\cite{Bacon2017OptionCritic} learns temporally extended actions with their own
termination conditions, a useful lens on the push primitive used here as a
macro-action. The survey of \cite{Hutsebaut2022HRLSurvey} situates these
approaches and their open challenges. Imagined-goal methods
\cite{Nair2018ImaginedGoals} generate goals from a learned latent, mirroring how
Alice's proposals define Bob's training distribution.

\subsection{Mechanics and Geometry of Planar Pushing}
\label{sec:pushing}

The final body of theory explains why the pushing task is genuinely
multi-objective. When a robot slides an object across a surface, the motion is
governed by friction between the object and the support. The analysis is made
tractable by the \textbf{quasi-static} assumption: pushing is slow enough that
inertial forces are negligible and the applied load is balanced entirely by
friction \cite{Mason1986Pushing}. Under this assumption the object moves only
while it is being pushed and stops the instant the push stops, so its trajectory
is determined by the geometry of frictional contact rather than by dynamics.

The central tool is the \textbf{limit surface} \cite{Goyal1991}. At a sliding
contact, the set of all friction forces and moments the interface can sustain
forms a convex surface in the three-dimensional load space of planar forces and
moment, $P = [F_x, F_y, M]$. For isotropic Coulomb friction this surface is a
convex, closed shape often approximated by an ellipsoid. Its importance lies in
the \textbf{normality (maximum-power) principle}: the object's instantaneous
motion twist $\mathbf{t} = [v_x, v_y, \omega]$, comprising a linear velocity and
an angular velocity, is normal to the limit surface at the point corresponding to
the applied load,
\begin{equation}
    \mathbf{t} \;\propto\; \nabla f(\mathbf{w}),
    \label{eq:normality}
\end{equation}
where $f(\mathbf{w}) = 0$ is the limit surface and $\mathbf{w}$ the applied wrench.
The consequence is important for this thesis. Consider any push not directed
exactly through the object's centre of friction. Its wrench lies where the surface
normal has a non-zero moment component, so the resulting twist contains an angular
velocity: \emph{almost every push rotates the object as well as translating it}.
Translation and rotation are mechanically coupled, and a goal that constrains both
cannot be decomposed into two independent sub-problems. This coupling is the
physical origin of the combined position-and-orientation success criterion that
recurs throughout the results. It is also the reason a reward or curriculum that
treats the two objectives as separable is inconsistent with the underlying
physics. \cite{LynchMason1996} developed the controllability and planning
consequences of this mechanics for stable pushing, and \cite{HoweCutkosky1996}
gave practical force-motion models. The survey of \cite{Stuber2020PushingSurvey}
and the high-fidelity dataset of \cite{Yu2016MillionWays} document how strongly
real pushing behaviour depends on object shape and pressure distribution.

This dependence on object shape distinguishes the two objects used in this thesis
and motivates one of its experiments. A rotationally symmetric disc has a radially
symmetric pressure distribution about its centre. A push whose line of action
passes through that centre applies a wrench with no moment component, so the
limit-surface normal points opposite the push direction and the resulting twist is
a pure translation. The disc therefore has no orientation to control, and a goal on
the disc constrains position alone. The T-block has an asymmetric mass
distribution, and few pushes pass through its centre of friction; each off-centre
push carries a moment, and by the normality condition of
Equation~\ref{eq:normality} the resulting twist contains an angular velocity.
Orientation is then an independent quantity that must be controlled alongside
position. This mechanical difference is the basis of the coupling-isolation
experiment of Chapter~\ref{section:methods}. The disc removes the orientation
objective at the level of the physics. A self-play model that reaches the
single-agent level on the disc, yet underperforms it on the T-block, thereby
identifies the coupled objective, rather than the object itself, as the source of
the difficulty.

Because position and orientation are coupled, it is natural to measure goal error
with a single metric on the space of planar poses, rather than with two separate
quantities. The configuration of a planar rigid body is an element of the
\textbf{Special Euclidean group} $\mathrm{SE}(2)$, the group of rigid-body
transformations of the plane. Its elements combine a translation $(x, y) \in
\mathbb{R}^2$ with a rotation $\theta \in \mathrm{SO}(2)$. \cite{Park1995LieGroup}
formulated rigid-body kinematics and dynamics in the language of Lie groups, which
provides the tools to define a distance on this space. A left-invariant Riemannian
metric on $\mathrm{SE}(2)$ yields a geodesic distance between two poses. Because
translation is measured in metres and rotation in radians, combining them requires
a \textbf{characteristic length} $L$ that converts an angular error into an
equivalent linear displacement. A pose error
can then be summarised by a single scalar of the form
\begin{equation}
    d_{\text{pose}} \;=\; \sqrt{\,\Delta x^2 + \Delta y^2 + L^2\, \Delta\theta^2\,},
    \label{eq:se2_distance}
\end{equation}
in which $L$ sets the relative weight of orientation against position and is
naturally identified with a physical length scale of the object. This unified
$\mathrm{SE}(2)$ distance underlies the single-potential reward used for the
self-play variants in Chapter~\ref{section:methods}. The idea of respecting the
geometry and symmetry of the configuration space also appears in recent
manipulation learning: SE(3)-equivariant and diffusion-based methods
\cite{Urain2023SE3Diffusion}, MDP homomorphic networks
\cite{VanderPol2020MDPHomomorphic}, equivariant transporter networks
\cite{Huang2022EquivariantTransporter}, and equivariant RL under partial
observability \cite{Nguyen2024EquivariantRL} all exploit the group structure of
the task, and provide broader context for a geometry-aware treatment of planar
pushing.

\subsection{Positioning of This Work}

The literature provides two well-founded but rarely compared techniques for
reducing the sample cost of on-policy manipulation learning: shaping the reward,
with the policy-invariance guarantee of PBRS, and structuring the training
process, through forced or automatic curricula such as ASP. Each has strong
theoretical support and prominent empirical results, yet the two are seldom
evaluated directly against each other on a single task, under a common budget,
and against a common success criterion. The planar-pushing setting makes the
comparison more stringent, because the limit-surface coupling of
Equation~\ref{eq:normality} makes the task irreducibly multi-objective and
therefore places a strong demand on any method that implicitly assumes the two
objectives can be handled separately. The chapters that follow implement both
techniques in this setting and measure what each contributes.
